\section{Validation and Testing Standards}
\label{sec:validation}
To make the validation story honest, we separate three levels of evidence:
\begin{itemize}
    \item \textbf{Code-path validation:} the training, generation, cleanup, export, and CFD scoring routines run end-to-end in this environment.
    \item \textbf{Solver validation:} the built-in CFD path can be compared against a higher-fidelity external solver when one is available.
    \item \textbf{Physical validation:} the generated shape should eventually be compared against analytic or experimental benchmarks, which is still future work.
\end{itemize}

\subsection{Verification}
Verification checks that the implemented steps are executed correctly. In this revision we verified that:
\begin{itemize}
    \item the shape-plausibility loss contributes to training,
    \item the connected-component cleanup hook reduces fragmented voxel outputs,
    \item STL export succeeds from the cleaned voxel grid, and
    \item the OpenFOAM export hook writes a runnable case directory skeleton, including `blockMeshDict`, `snappyHexMeshDict`, and surface STL output.
\end{itemize}

\subsection{Validation}
Validation is partial but no longer purely blocked. The built-in CFD path was exercised on the reduced CPU pipeline, and a reduced OpenFOAM sonicFoam benchmark was run on the corresponding geometry. The comparison used the lower-level `forces` output and then normalized the result manually into coefficients. That makes the comparison approximate, but it is still useful as a solver-to-solver sanity check.

The paper therefore claims the following and nothing more: the export path is concrete, the benchmark comparison is measurable, and the current internal solver diverges substantially from the OpenFOAM pressure estimate on the reduced cube test. By being explicit about the manual extraction method and the remaining convention differences, the revision avoids overstating the credibility of the aerodynamic numbers.

